% Generated from project outputs. Do not edit by hand.
\begin{table}[!htbp]
\centering
\begin{threeparttable}
\caption{Labor income by subsample: 2024 moments}
\label{tab:annual-2024-y1-moments}
\small
\setlength{\tabcolsep}{4pt}
\renewcommand{\arraystretch}{1.15}
\begin{tabularx}{\linewidth}{@{}Xrrrrr@{}}
\toprule
Sample & N & Mean & SD & Min & Max \\
\midrule
General & 92 & 3,288.1 & 1,695.0 & 746.7 & 7,396.2 \\
Men & 92 & 3,805.2 & 1,710.4 & 841.9 & 8,511.5 \\
Women & 92 & 2,629.3 & 1,733.3 & 482.6 & 7,207.4 \\
Urban & 92 & 3,499.3 & 1,678.0 & 934.6 & 6,764.5 \\
Rural & 92 & 2,848.0 & 1,800.5 & 525.0 & 14,014.4 \\
Indigenous & 92 & 3,288.1 & 1,695.0 & 746.7 & 7,396.2 \\
Bottom 99 percent & 92 & 3,080.9 & 1,428.1 & 746.7 & 5,861.9 \\
Bottom 90 percent & 92 & 2,540.5 & 903.7 & 746.7 & 4,422.5 \\
Bottom 50 percent & 92 & 1,363.2 & 372.3 & 398.9 & 2,172.3 \\
\bottomrule
\end{tabularx}
\begin{tablenotes}[flushleft]
\footnotesize
\item Monthly quetzales. Unweighted distribution of survey-weighted cohort means before transition matching. N counts cohort cells, not individuals. SD is dispersion, not a standard error. Subsamples overlap. Cumulative income definitions and treatment of missing components follow the merging scripts.
\end{tablenotes}
\end{threeparttable}
\end{table}
